
\section{Closing a year without erasing a missed payment}
\label{sim:financing-carry}
A reporting boundary creates another opportunity to confuse present balances
with financial history. Return to the late-funded borrower. In the paid case,
the final settlement leaves no principal outstanding. In the unpaid case,
the company has received the new equity but retains cash 40, outstanding
principal 40 and arrears 40 at year-end. Both cases record new shares of 40
and a historical missed payment of 40. The second case must be able to enter
the next year with the unpaid obligation intact; requiring every claim to be
settled before carrying the book would exclude the scenario being studied.

Closing transfers interest expense of 20 into retained earnings under the
selected report mapping. Current-period flows are archived before the new
period begins. Carry then initializes the next reporting period without
turning that accounting boundary into a payment, waiver or new issue. The
selected policy preserves outstanding principal, arrears, share quantities
and the old missed amount. It also preserves the prior clause category,
observation date, observed economic-state identifier, supplied evidence
revision and acceleration record.

This policy is explicit because a generic typed carry operation may permit
different treatment of different memory fields. A type-correct reset of the
missed amount to zero would destroy financially relevant history while leaving
the opening balances coherent. The selected financing wrapper therefore
requires preservation for these fields before admitting the next opening.
Other memory, including period-specific counters, can require a different
policy; preservation of this selected contract record is not a universal
fiscal-closing rule.

Preserving the observation does not make it current. The new period and its
flow state change the economic dependency in
equation~\eqref{sim:remedy-observation-dependency}, and the old observation date
remains old. A new proposed acceleration consequently requires a new
assessment. Even after that assessment, a changed supplied evidence revision
requires reassessment again. Carry must not refresh these fields merely to
make a later command pass.

The unpaid example continues into 2028 with an explicit fresh observation
and a payment of 40. Payment reduces cash and the outstanding claim together;
the historical miss remains 40. The separately recorded acceleration also
survives unless an applicable rescission or other supported rights-changing
event changes it. A principal-payment observation still cannot authorize the
automatic bankruptcy route. Those are different contractual premises even
when a reporting boundary lies between their observations.

The exact continuation compares the complete closing entries, archived flows,
next opening, subsequent payment and next-period statements on two account
charts. The unchanged carry wrapper has separate compatibility evidence for
the stronger remedy profile that includes supplied evidence revision and
clause category. This is explicit compatibility for the selected records,
not an implicit migration between earlier and later schemas. The examples
demonstrate why financial histories and observation dependencies must travel
together across reporting periods. Actual terms, several simultaneous
clauses, guarantees and the economic consequences of acceleration still
require their own sources and methods.


\section{Conversion changes a claim and the people who control it}
\label{sim:conversion-book-rights}
Preserving rights across a reporting boundary is necessary until an event
changes them. A conversion provides a contrasting case: cash may stay fixed
while debt disappears, new shares are issued and another holder gains
control. A simulator that checks cash and balanced journals alone can miss
the decision consequence.

Consider a separately supplied fiscal opening on 1 January 2028. Project
assets are 100 and cash is zero. Senior debt has face and carrying amount
40. A convertible claim has face 80 and a supplied carrying amount $C$.
There are 40 old shares and no shares held by the convertible creditor.
Opening capital is $60-C$, so the signed balances reconcile. The examples
use $C=80$ and $C=17000/441$ as two distinct accounting premises. The latter
is a diagnostic amount; selecting it does not establish a measurement rule.

The supplied contract says that the conversion signal extinguishes the
claim and issues 120 shares to its holder on that date. Under the selected
carrying-value transfer, the journal debits the convertible liability by
$C$ and credits issued equity by $C$. The face becomes zero, new shares
become 120 and the observed signal and conversion are recorded. There is
no cash flow or profit from this assumed transfer. The financial statements
then show assets 100, liabilities 40 and equity 60 in either carrying case.
Under the other signal, the debt remains outstanding and no shares issue.

Cash-flow and voting fractions follow from the resulting shares under the
assumed equal rights per share. After conversion the new holder owns
$120/(40+120)=3/4$ and the old holder owns $1/4$. Under this supplied economic governance rule, the new holder therefore
has strict majority control. Before conversion the old holder controls all
shares. The separate accounting-control assessment, including potential
rights before legal conversion, is explained in
Section~\ref{sim:conversion-control-scope}. These rights follow from the completed transition, rather than
from a controller label retained from the opening. Different voting classes
would require a different supplied rights rule.

The complete comparisons retain every account, the claim and share record,
unrelated memory, event history, journal and report on two account charts.
They also compare the resulting rights with a separately authored payoff
table. Omitting claim extinction or share issue fails even when a proposed
journal balances. Repeating conversion with a fresh event identifier fails
because the event has already occurred. Wrong dates and signals fail before
effects are accepted. Reporting the equity increase without explaining its
noncash movement also fails, while preserving the already accepted conversion.

This example starts from an explicit opening and assumes the accounting
transfer. It does not derive that opening from issuance, select a reporting
standard, or establish that a real contract has these terms. Its contribution
is narrower and useful: the accounting balances and economic rights change
together, and the rights can supply the next decision without counting both
converted debt and equity as surviving claims.

